\section{How Ratio Compares}

The incumbents win on breadth today --- years of integrations, custodian feeds,
and reporting templates. Ratio wins on the foundation, and on the dimensions that
breadth cannot retrofit: verifiable correctness, reproducibility, openness, and
economics. The table below is a capability view, not a feature checklist;
it asks ``can the platform let you \emph{verify} this,'' not ``does it have a
button for it.''

% Capability comparison matrix. \yes full, \partialmark partial, \no none.
\begingroup
\renewcommand{\arraystretch}{1.35}
\setlength{\tabcolsep}{5pt}
\begin{center}
\sffamily\footnotesize
\begin{tabular}{@{}p{0.40\textwidth}*{5}{>{\centering\arraybackslash}p{0.085\textwidth}}@{}}
\rowcolor{ratioInk}
\textbf{\color{white}Capability} &
\textbf{\color{white}Ratio} &
\textbf{\color{white}Black Diamond} &
\textbf{\color{white}Orion} &
\textbf{\color{white}Addepar} &
\textbf{\color{white}Tamarac} \\
\arrayrulecolor{ratioSand}\midrule
Investment accounting: multi-currency, positions, tax lots, P\&L &
\yes & \yes & \yes & \yes & \yes \\
\rowcolor{ratioCream}
Performance reporting \& client portal &
\partialmark & \yes & \yes & \yes & \yes \\
Formally verified (machine-checked) accounting kernel &
\yes & \no & \no & \no & \no \\
\rowcolor{ratioCream}
Bit-exact, reproducible historical replay &
\yes & \no & \no & \partialmark & \no \\
Content-addressed, immutable audit trail &
\yes & \partialmark & \partialmark & \partialmark & \partialmark \\
\rowcolor{ratioCream}
Open core, self-hostable, no data lock-in &
\yes & \no & \no & \no & \no \\
Inspectable engine (read the arithmetic) &
\yes & \no & \no & \no & \no \\
\rowcolor{ratioCream}
AI authoring with the model fenced from the ledger &
\yes & \partialmark & \partialmark & \partialmark & \no \\
Extensible in code (Rust / Python) &
\yes & \no & \partialmark & \partialmark & \no \\
\rowcolor{ratioCream}
Pricing decoupled from assets under management &
\yes & \no & \no & \no & \no \\
\arrayrulecolor{ratioSand}\bottomrule
\end{tabular}

\smallskip
{\color{ratioBrown}\footnotesize \yes\ full \quad \partialmark\ partial / emerging \quad \no\ not offered (or not inspectable)}
\end{center}
\endgroup


\subsection{Where the categories sit}
Plotted against the two axes that matter most for the next decade --- how
\emph{open and inspectable} a platform is, and how much \emph{provable assurance}
it offers --- the incumbents cluster in the established, closed-but-capable
quadrant. Ratio occupies the one quadrant none of them can reach without
rebuilding their engine.

% Positioning quadrant — Ratio's view. x: openness, y: provable assurance.
\begin{figure}[h]
\centering
\begin{tikzpicture}[
  >={Stealth[length=2.2mm]}, font=\sffamily\footnotesize,
  dot/.style={circle, fill=ratioBrown, inner sep=2pt},
  lbl/.style={font=\sffamily\footnotesize, text=ratioInk},
]
% Desirable quadrant shading (top-right)
\fill[ratioSand!45, rounded corners=2pt] (3,3) rectangle (6,6);
\node[align=center, text=ratioBrown, font=\sffamily\bfseries\small] at (4.5,5.55)
  {Verifiable \& open\\[-1pt]\footnotesize\mdseries (the next decade)};

% Axes
\draw[->, ratioInk, line width=0.9pt] (0,0) -- (6.4,0)
  node[right, lbl, align=center] {};
\draw[->, ratioInk, line width=0.9pt] (0,0) -- (0,6.4);
\node[lbl, align=center] at (3.2,-0.55) {Openness \& inspectability $\rightarrow$};
\node[lbl, align=center, rotate=90] at (-0.55,3.2) {Provable assurance \& reproducibility $\rightarrow$};

% Incumbent cluster (lower-left / center)
\node[dot] (bd) at (1.5,1.7) {}; \node[lbl, right=1mm of bd] {Black Diamond};
\node[dot] (or) at (1.9,1.2) {}; \node[lbl, right=1mm of or] {Orion};
\node[dot] (ta) at (1.2,1.1) {}; \node[lbl, left=1mm of ta] {Tamarac};
\node[dot] (ad) at (2.4,2.6) {}; \node[lbl, right=1mm of ad] {Addepar};
\node[dot, fill=negGray] (ss) at (1.05,0.65) {}; \node[lbl, anchor=west, text=negGray] at (1.2,0.5) {spreadsheets / in-house};

% Ratio — emphasized, alone, top-right
\node[circle, fill=posGreen, inner sep=3.4pt, label={[text=ratioInk,font=\sffamily\bfseries]right:Ratio}]
  (rt) at (4.6,4.7) {};
\draw[->, posGreen, line width=0.9pt, densely dashed] (ad) to[bend left=15] (rt);
\end{tikzpicture}
\caption*{\small\color{ratioBrown}\sffamily Ratio's positioning view. Incumbents
are capable but closed; only Ratio combines an open, inspectable engine with
machine-checked assurance and reproducible history.}
\end{figure}


\begin{callout}
\footnotesize\color{ratioInk}
\textbf{A note on fairness.} Comparisons reflect Ratio's interpretation of
publicly available information as of 2026 and general patterns among legacy
wealth-technology platforms; capabilities evolve and specific deployments vary.
Black Diamond is a trademark of SS\&C Technologies; Orion, Addepar, and Tamarac
are trademarks of their respective owners. Ratio is an independent project and is
not affiliated with, sponsored by, or endorsed by any of them. Where a competitor
offers a partial capability, we mark it \partialmark\ rather than \no.
\end{callout}
